\section{Late information: learning deadlines and their exact price}\label{sec:deadline}

Observation is a design variable. In the delayed instance the regime is revealed exogenously; in many sustainability settings---an unknown regeneration rate, an uncertain climate sensitivity, an institution whose effectiveness is unknown---the hidden parameter can only be learned by acting, and acting on a wrong hypothesis is costly. This section prices that structure exactly on a second instance and then states the general threshold law behind it.

\subsection{The declining instance}

A stock $z$ drifts down by $\tfr1{10}$ per stage unless corrected; the correction is gated by a hidden sign $\theta\in\{-1,+1\}$:
\begin{equation}\label{eq:declining}
z_{t+1}=z_t-\tfr1{10}+\theta\,u_t,\qquad u_t\in\{-1,0,+1\},\qquad\text{safe set }\{z\ge1\}.
\end{equation}
Holding $u=0$ is blind: the displacement $-\tfr1{10}$ is the same on both branches. Any $u\neq0$ is a \emph{probe}: the displacement $-\tfr1{10}\pm1$ separates the branches by $2$, so a single reading of the displacement reveals $\theta$ exactly and the posterior collapses---one observation splits the belief with no statistics. The \emph{learned} control $u=\theta$ has drift $-\tfr1{10}+\theta^{2}=\tfr9{10}>0$ and is safe forever from any $z\ge1$; the wrong control $u=-\theta$ has drift $-\tfr{11}{10}$.

\begin{proposition}[probes and learning deadlines]\label{prop:deadline}
On the instance \eqref{eq:declining} with the grid $z_0\in\{1.0,\dots,2.5\}$:
\begin{enumerate}[leftmargin=1.6em,itemsep=1pt]
\item[(a)] with $\theta$ observed from the start every $z_0\ge1$ is viable ($16$ of $16$ grid states);
\item[(b)] with a free probe (probe at stage $0$, then the learned control) viability holds exactly when $z_0\ge\tfr{21}{10}$, independently of any deadline ($5$ grid states);
\item[(c)] in the \emph{pinned-probe family}---hold $u=0$ to the deadline $\Tl$, probe at $\Tl$, then the learned control, no earlier probe being admissible---viability holds exactly when
\begin{equation}\label{eq:deadline}
z_0\ \ge\ \tfr{21}{10}+\tfr{\Tl}{10},\qquad \Tl=0,1,2,\dots,
\end{equation}
so the viable set on the grid shrinks linearly and antitonically in the deadline, $|K|=5,4,3,2,1$ for $\Tl=0,\dots,4$;
\item[(d)] with probe-free exogenous revelation at $\Tl$ (hold $u=0$, then the learned control) viability holds exactly when $z_0\ge1+\tfr{\Tl}{10}$, with $|K|=16,15,14,13,12$ for $\Tl=0,\dots,4$.
\end{enumerate}
The pinned family sits between the two benchmarks: the flat free-probe threshold prices the probe excursion alone, the probe-free threshold prices the blind decline alone, and \eqref{eq:deadline} prices both (Figure~\ref{fig:deadlines}, Table~S4.1).
\end{proposition}

\begin{proof}
(a) $u=\theta$ from $z_0\ge1$ never exits. (b) The probe moves the branches to $z_0-\tfr1{10}\pm1$; the $\theta=-1$ branch requires $z_0-\tfr{11}{10}\ge1$. (c) Before the deadline the state falls to $z_0-\tfr{\Tl}{10}$, which must stay at or above the floor; the probe then requires $z_0-\tfr{\Tl}{10}-\tfr{11}{10}\ge1$, the binding condition; afterwards the learned control is safe. The branch positions after the probe differ by $2$, so the reading identifies $\theta$. The grid counts follow from \eqref{eq:deadline}. (d) Only the blind decline must be survived: $z_0-\tfr{\Tl}{10}\ge1$.
\end{proof}

\begin{figure}[t]
\centering
\includegraphics[width=\linewidth]{fig3_deadlines.pdf}
\caption{Learning deadlines on the declining instance. (a) The pinned-probe threshold $\tfr{21}{10}+\tfr{\Tl}{10}$ of Proposition~\ref{prop:deadline} between the flat free-probe benchmark $\tfr{21}{10}$ and the probe-free benchmark $1+\tfr{\Tl}{10}$. (b) Viable grid states: $5,4,3,2,1$ for the pinned probe and $16,\dots,12$ for exogenous revelation at $\Tl=0,\dots,4$; beyond $\Tl=4$ the pinned family is empty on the grid. The grid counts form a staircase, not a continuous law; the thresholds are exact for every real $z_0$.}
\label{fig:deadlines}
\end{figure}

\subsection{The general schedule law}

\begin{theorem}[exact threshold of a hold--probe--learn schedule]\label{thm:additive}
Let the hidden parameter $\theta$ range over a finite set $\Theta$ and enter the drift additively, $z_{t+1}=z_t+d(u_t,\theta)$, with floor $z\ge1$. Consider the schedule: hold $u_0$ for $T\ge0$ stages, probe with $u_p$ for one stage (the displacement $d(u_p,\cdot)$ being injective on $\Theta$, so the probe reveals $\theta$), then play $u^{*}(\theta)$ with $\mu=\min_{\theta}d(u^{*}(\theta),\theta)>0$. Then the schedule is viable for every horizon if and only if
\begin{equation}\label{eq:general}
z_0\ \ge\ 1+\max_{\theta\in\Theta}\max\bigl\{0,\ -T\,d(u_0,\theta),\ -T\,d(u_0,\theta)-d(u_p,\theta)\bigr\}.
\end{equation}
Let $d_0=\max_\theta(-d(u_0,\theta))\ge0$ be the worst hold loss and $e_p=\max_\theta(-d(u_p,\theta))\ge0$ the worst probe excursion. If some $\theta^{*}$ attains both maxima (\emph{aligned worst losses}), \eqref{eq:general} is the additive law
\begin{equation}\label{eq:additive}
z_0\ \ge\ 1+T\,d_0+e_p .
\end{equation}
Without alignment the additive right-hand side is only an upper bound on the exact threshold, and the gap can be positive.
\end{theorem}

\begin{proof}
On each branch $\theta$ the path is linear within each phase, so its minimum over the hold--probe prefix is attained at a phase endpoint: at $z_0$, at $z_0+T\,d(u_0,\theta)$, or at $z_0+T\,d(u_0,\theta)+d(u_p,\theta)$; after the probe the learned drift $\mu>0$ lifts the state away from the floor forever. Viability on every branch is therefore the conjunction of the three endpoint conditions over $\theta$, which is \eqref{eq:general}. Under alignment the third term at $\theta^{*}$ equals $T d_0+e_p$, which dominates the other terms because $d_0,e_p\ge0$, giving \eqref{eq:additive}; in general $\max_\theta(-Td(u_0,\theta)-d(u_p,\theta))\le Td_0+e_p$, with strict inequality possible.
\end{proof}

\begin{example}[two aligned instances and a misaligned one]\label{ex:additive}
The declining instance has $d_0=\tfr1{10}$ (attained at both $\theta$) and $e_p=\tfr{11}{10}$ (attained at $\theta=-1$), so the losses are aligned and \eqref{eq:additive} is \eqref{eq:deadline}. A second instance, $z_{t+1}=z_t-\tfr3{20}+\tfr9{10}\,\theta u_t$, has $d_0=\tfr3{20}$ and $e_p=\tfr{21}{20}$ and threshold $\tfr{41}{20}+\tfr{3T}{20}$; both are verified by exact enumeration for all swept $(z_0,T)$. Alignment is essential: with $\Theta=\{-1,+1\}$, hold drift $d(u_0,\theta)=-\tfr12+\tfr12\theta$, probe drift $d(u_p,\theta)=-\tfr12-\tfr12\theta$ and learned drift $1$, the worst hold loss $d_0=1$ is attained at $\theta=-1$ and the worst probe excursion $e_p=1$ at $\theta=+1$; the additive right-hand side at $T=1$ is $3$, but \eqref{eq:general} gives the exact threshold $2$, and indeed from $z_0=2$ both branches follow the paths $(2,1,1)$ and $(2,2,1)$ and then rise.
\end{example}

\begin{remark}[the delayed instance as a schedule]\label{rem:delayedschedule}
The delayed instance of Section~\ref{sec:delayed} under the hold class is the schedule of Theorem~\ref{thm:additive} with $T=\Tobs$, no probe ($e_p=0$), hold loss $d_0=1$ on the falling branch and the learned control $u=\theta$ afterwards: \eqref{eq:additive} gives $z_0\ge1+\Tobs$, the locus of the timing certificate and of the value-one boundary of \eqref{eq:hold} at $k>\Tobs$. The same law therefore prices exogenous and endogenous revelation on one scale: a unit of blind decline per stage, plus the excursion of whatever experiment produces the information.
\end{remark}
